% Содержимое отчета по курсу Анализ алгоритмов

\aaunnumberedsection{ВВЕДЕНИЕ}{sec:intro}

Многопоточность — это метод параллельного выполнения, при котором несколько потоков обрабатывают
задачи одновременно в рамках одного процесса, используя общую память и ресурсы.
Она позволяет ускорить выполнение задач, которые могут выполняться параллельно,
таких как обработка больших объемов данных и выполнение сложных вычислений.
Благодаря многопоточности приложения могут эффективно задействовать все ядра процессора,
что особенно важно для многопроцессорных систем~\cite{paralel_cpp}.

Нативные потоки (Native Threads) — это потоки, управляемые непосредственно операционной системой.
Использование нативных потоков позволяет создавать высокопроизводительные многопоточные приложения,
но требует контроля доступа к общим ресурсам для предотвращения гонок данных.

Цель работы --- получение навыка организации параллельных вычислений на основе нативных потоков.
Для достижения поставленной цели необходимо выполнить следующие задачи:
\begin{itemize}
    \item[---] описать входные и выходные данные;
    \item[---] реализовать два алгоритма для загрузки контента из html--страниц -- последовательный и параллельный с использованием нативных потоков;
    \item[---] протестировать разработанные алгоритмы;
    \item[---] сравнить скорость выполнения программы в зависимости от количества используемых потоков.
\end{itemize}

\aasection{Входные и выходные данные}{sec:input-output}

Входными данными программы являются адрес главной страницы ресурса~\cite{recepies}, а также количество потоков и максимальное количество страниц, с которых выгружаются данные.
Выходными данными является директория, содержащая скачанные данные со страниц в формате, пригодном для дальнейшей обработки (html)

\aasection{Преобразование входных данных в выходные}{sec:algorithm}

Потоки выполняют скачивание очередной страницы, сохранение результата в файл,
а также поиск на странице ссылок на другие страницы с добавлением ссылок в структуру данных типа множество строк.

\aasection{Тестирование}{sec:tests}

В таблице~\ref{tbl:tests} представлены функциональные тесты для разработанного ПО.
Все тесты пройдены успешно.

\begin{table}[h]
    \begin{center}
        \caption{\label{tbl:tests} Функциональные тесты}
        \begin{tabular}{|p{5cm}|p{5cm}|p{5cm}|}
            \hline
            Описание~теста & Выходные~данные & Ожидаемые~выходные~данные
            \\ \hline
            Число страниц~для~загрузки <= 0 & Сообщение~об~ошибке & Сообщение~об~ошибке
            \\ \hline
            Число~потоков <= 0 & Сообщение~об~ошибке & Сообщение~об~ошибке
            \\ \hline
            Некорректный~url первой~страницы & Сообщение~об~ошибке & Сообщение~об~ошибке
            \\ \hline
            Работа~в~параллельном~режиме, 4~потока & Директория~$data$~c 4 файлами & Директория~$data$~c 4 файлами
            \\ \hline
            Работа~в~последовательном~режиме, 1~поток & Директория~$data$~c 1 файлом & Директория~$data$~c 1 файлом
            \\ \hline
        \end{tabular}
    \end{center}
\end{table}

%\begin{table}[h]
%    \begin{center}
%        \caption{\label{tbl:tests} Функциональные тесты}
%        \begin{tabular}{|l|l|l|}
%            \\\hline
%            Описание~теста & Выходные~данные & Ожидаемые~выходные~данные
%            \\\hline
%            Число страниц для загрузки <= 0 & Сообщение об ошибке & Сообщение об ошибке \\\hline
%            Число потоков <= 0 & Сообщение об ошибке & Сообщение об ошибке & \\\hline
%            Некорректный url первой страницы & Сообщение об ошибке & Сообщение об ошибке \\\hline
%            Работа в параллельном режиме, 4 потока & Директория $data$ c 4 файлами & Директория $data$ c 4 файлами \\\hline
%            Работа в последовательном режиме, 1 поток & Директория $data$ c 1 файлом & Директория $data$ c 1 файлом \\\hline
%        \end{tabular}
%    \end{center}
%\end{table}

%\begin{longtable}{|p{.2\textwidth - 2\tabcolsep}|p{.33\textwidth - 2\tabcolsep}|p{.24\textwidth - 2\tabcolsep}|p{.23\textwidth - 2\tabcolsep}|}
%    \caption{Функциональные тесты}\label{tbl:tests} \\\hline
%    Описание~теста & Выходные~данные & Ожидаемые~выходные~данные \\\hline
%    \endfirsthead
%    \endfoot
%    Число страниц для загрузки <= 0 & Сообщение об ошибке & Сообщение об ошибке \\\hline
%    Число потоков <= 0 & Сообщение об ошибке & Сообщение об ошибке & \\\hline
%    Некорректный url первой страницы & Сообщение об ошибке & Сообщение об ошибке \\\hline
%    Работа в параллельном режиме, 4 потока & Директория $data$ c 4 файлами & Директория $data$ c 4 файлами \\\hline
%    Работа в последовательном режиме, 1 поток & Директория $data$ c 1 файлом & Директория $data$ c 1 файлом \\\hline
%\end{longtable}

\aasection{Описание исследования}{sec:study}

В ходе исследования требуется сравнить скорость выполнения программы в зависимости от количества используемых потоков.
Замеры проводились при обработки 100 страниц. Число дополнительных потоков изменялось с 1 до 16.
По результатам исследования был построен график с использование библиотеки $matplotlib python$~\cite{plot}\newpage


\section{Технические характеристики}
Технические характеристики используемого устройства:
\begin{itemize}
    \item процессор: AMD Ryzen™ 7--5800H with Radeon™ Graphics × 16;
    \item ОС: 64-разрядная операционная система Ubuntu 24.04.1 LTS;
    \item оперативная память: 16,0 ГБ;
    \item число логических ядер -- 16.
\end{itemize}

\section{Результаты исследования}

В таблице~\ref{tbl:research} приведено время выполнения
программного обеспечения в секундах.

\begin{table}[h]
    \begin{center}
        \caption{\label{tbl:research} Сравнение скорости работы алгоритма от числа потоков}
        \begin{tabular}{|c|r|}
            \hline
            Число потоков & Время, секунды
            \\ \hline
            1 & 18.42
            \\ \hline
            2 & 11.87
            \\ \hline
            4 & 6.07
            \\ \hline
            8 & 3.37
            \\ \hline
            12 & 2.52
            \\ \hline
            16 & 2.32
            \\ \hline
        \end{tabular}
    \end{center}
\end{table}

На рисунке~\ref{img:threads} показана зависимость времени работы от количества
потоков с постоянным количеством обрабатываемых страниц.

\FloatBarrier
\includeimage
{threads}
{f}
{h}
{\textwidth}
{Зависимость времени работы от количества потоков.}
\FloatBarrier

\aaunnumberedsection{ЗАКЛЮЧЕНИЕ}{sec:end}

В ходе исследования было установлено, что многопоточная реализация работает гораздо быстрее последовательной
(примерно в 4,5 раза), скорость обработки страниц возрастает в 1.5 раза при увеличении числа потоков в 2 раза.

Цель работы была достигнута.
Решены все поставленные задачи:
\begin{itemize}
    \item описаны входные и выходные данные;
    \item реализованы 2 алгоритма обработки страниц: многопоточный и последовательный;
    \item разработанные алгоритмы протестированы;
    \item выполнено сравнение скорости выполнения программы в зависимости от количества используемых потоков.
\end{itemize}
